\subsection*{64. The Bowling-Alley Intuition and Its Limits}

Geoffrey Moore's bowling-alley metaphor imagines a successful initial niche creating momentum into nearby niches [R1]. A credible reference in 1 regulated team may persuade similar teams; a specialized product can be extended to the next related workflow.

But bowling pins do not actually fall because they have been drawn adjacent on a slide. A hospital administrator may not trust references from a software startup. A beginner student does not necessarily value an advanced AI engineer's success story. A college procurement committee may require approvals that a D2C buyer does not.

For each proposed next pin, ask: \textbf{What specific asset produced by the first market reduces a purchase or delivery barrier in the second?} If you cannot answer, the next market is a fresh startup with some shared overhead.

\subsection*{65. A Transfer Matrix for Adjacent Opportunities}

Evaluate potential adjacencies across 8 core dimensions:

\begin{table*}[t]
\centering
\small
\begin{tabularx}{\textwidth}{l p{6.5cm} X}
\toprule
\textbf{Transfer Dimension} & \textbf{Questions to Investigate} & \textbf{High Transfer Looks Like} \\
\midrule
Customer relationship & Is the same payer or buyer group involved? & Existing satisfied account can fund the new job \\
Problem and outcome & Is the same core competence needed? & Existing solution needs moderate configuration only \\
Product and IP & Which assets are directly reusable? & Rubrics, software modules, processes, materials survive \\
Distribution & Can the same channel reach buyers? & Referrals and partners introduce new prospects naturally \\
Evidence and brand & Does existing proof persuade this buyer? & References are viewed as relevant to the new stakes \\
Delivery & Can the current team execute? & No new highly specialized permanent role required \\
Economics & Are sales and support patterns similar? & Comparable contribution margins and cash cycles \\
Governance & Do new compliance obligations arise? & Legal, security, and contractual terms remain manageable \\
\bottomrule
\end{tabularx}
\caption{Systematic transfer matrix for evaluating candidate adjacencies.}
\label{tab:ch7-transfer-matrix}
\end{table*}

A single decisive incompatibility can outweigh 6 superficial similarities. The matrix exposes the true cost of expansion.

\subsection*{66. When to Add the Next Segment: 4 Economic Signals}

A practical expansion decision benefits from 4 signals:

\begin{enumerate}
    \item \textbf{The core has repeatable contribution:} Future expansion is not financed by sacrificing core delivery or masking ongoing losses.
    \item \textbf{The core has spare or expandable capacity:} The founder can delegate or run a reliable operating system, rather than adding all new work personally.
    \item \textbf{The adjacent opportunity has paid evidence:} Real buyer commitments, or a bounded credible path to obtaining them, not merely compliments.
    \item \textbf{The expansion has a risk budget:} Enough cash and time are ring-fenced for the experiment without risking fulfillment, payroll, refunds, or essential reserves.
\end{enumerate}

There is no universal revenue milestone at which a second offering becomes wise. A company at ₹30 lakh annual revenue might responsibly test a tiny adjacent workshop. A ₹30 crore business can still be too fragile to expand if it depends on 1 client and consistently misses delivery obligations.

\subsection*{67. Expansion Readiness Is a Vector, Not a Single Score}

One business may have excellent product repeatability but weak acquisition. Another may have cash reserves but no management bench. A third may have strong demand but inadequate instructor capacity. Summing these into a ``growth readiness score of 82'' hides the binding constraint.

Use a vector: \textbf{product proof, channel repeatability, contribution, liquidity, delivery capacity, management capacity, risk controls, and adjacency evidence.} Record each separately and define minimum acceptable conditions for the proposed expansion.

The desired threshold depends on the irreversibility of the commitment. A 4-week pilot may require modest proof. Signing a multiyear lease, hiring a large sales team, opening a campus, or guaranteeing training capacity to a university requires far stronger evidence and reserves.

\subsection*{68. When Not to Expand}

Do not launch the second business simply because the founder is bored with the first. Do not add products to compensate for weak product-market fit. Do not enter the nearest apparent adjacency without studying its buyer, budget, delivery, and economics. Do not confuse ``we can teach this'' with ``a qualified segment will buy this at a sustainable price.''

Other red flags: the core requires weekly founder rescue; cash collections are uncertain; the new offering depends on the same overcommitted experts; customers routinely report disappointing outcomes; or all expansion revenue is expected from a single friendly buyer. Under these conditions, growth can magnify weaknesses faster than it adds assets.

\subsection*{69. When It Is Rational to Change the Beachhead}

A pivot is rational when evidence undermines a decisive premise: the trigger is too rare, qualified buyers cannot pay, the route to them is structurally closed, delivery economics fail even after sensible redesign, or the market needs a different mechanism than the founder can responsibly provide.

A pivot is not an admission that the original ambition was foolish. It is evidence that 1 pathway is weak. Preserve transferable learning: interview histories, technical prototypes, customer relationships, delivery tools, and negative findings. Record the rationale to prevent later nostalgia from rewriting history.

Do not change segments after every rejection. Use experiments large or deep enough to investigate the specific uncertainty, and compare results against explicit stop or redesign conditions established before the outcome was known.

\subsection*{70. The Long-Term Paradox: Become Known for Something Before Everything}

At the beginning, customers may only recognize an institution in connection with 1 narrowly defined result. Later, after repeated delivery and accumulated proof, its name can extend to related capabilities. Eventually a mature institution may span many disciplines or research fields while maintaining coherence through culture, governance, and standards.

The order matters: \textbf{Breadth is easier to finance and believe after depth has generated evidence, capability, and cash.} However, the destination need not be endless diversification. Some exceptional businesses remain narrow specialists for decades and create substantial wealth and impact.

The founder's question is not how quickly to stop being focused. It is whether the next expansion inherits enough from the last one to become stronger rather than merely larger.
